% How the model translates the C code and how the contracts speak about it; then the parts of the
% program around the sieve: the command line, the curve searched, the text of a point, the exact
% check of a triple, and the program with find_points.
\chapter{The model and its contracts}\label{chap:model}

This chapter describes how the model translates \jpoints{} 3.0.1 into Lean and how the proofs
speak about it, and then the parts of the program that lie around the sieve: the reading of the
command line, the curve that the search runs on, the text of a point, the exact check of a triple,
and the program with its search. The sieve itself, whose task is to drop triples over which no
point lies, is the subject of Chapters~\ref{chap:sieve-numbers} to~\ref{chap:region}; the record of
which declaration models which lines of the C code, and the tests of the model against the
program, are in Chapter~\ref{chap:ccode}.

\section{How the C code is modeled}\label{sec:modeling}

\paragraph{Plain code, function by function.} The model is plain Lean on core and \texttt{Std}: no
Mathlib, and no \texttt{partial}, \texttt{unsafe}, \texttt{implemented\_by} or \texttt{extern}, so
that the executable \texttt{j-points-model} runs exactly the code that the proofs see. Each C
function is a Lean definition, or several where it is split at its loops and its switches; each
docstring names the lines of the C code that it models and says, in a paragraph headed
\emph{Deviation}, where and why the model departs from them (Chapter~\ref{chap:ccode}). The globals
of \jpoints{} are the fields of structures: the command line (\lname{JPoints.Cli.Input}), the curve
searched (\lname{JPoints.SearchCurve}), what the exact check reads (\lname{JPoints.CheckGlobals}),
the plan of the sieve (\lname{JPoints.Plan}), the state of the thread (\lname{JPoints.SiftState})
and others. They are computed in the order of the C code and passed on, and a function returns
what it writes. Pointers become indices, the text that the C code prints on the standard output is
collected in an output (\lname{JPoints.Output}), and an exit through \texttt{error} is a value
returned.

\paragraph{Components in \texttt{Id}.} A component with a loop is a definition of
type~$\mathtt{Id}\ \alpha$ in Lean's \texttt{do} notation, with mutable variables: a \texttt{for}
loop over a range where the C code has a counted loop, and a \texttt{while} loop where its loop
ends on the data, whose end the contracts prove by a measure. Its callers bind its result; the
program runs the whole (Definition~\ref{def:Program.run}), and the examples beside the model run
single components and compare their values with those of the C code.

\paragraph{Contracts.} The proofs, on Mathlib, give each component a contract: a Hoare triple
$\{P\}\ c\ \{Q\}$ of \texttt{Std.WP}, marked \texttt{@[spec]}, which says that when the
precondition~$P$ holds, the result~$r$ of the component~$c$ satisfies~$Q\,r$. It is proved by the
tactic \texttt{vcgen}, which unfolds the component, replaces each call by the contract of the
callee, and reduces the triple to verification conditions. A loop needs an invariant, a named
definition on the tuple of the mutable variables that the loop assigns, and a \texttt{while} loop
also a measure that decreases. A tactic of the project, \texttt{vc\_simp}, rewrites the encodings
of the weakest preconditions in the conditions into plain propositions, and each condition is
discharged by a lemma stated in its own shape, so that the proofs depend neither on the order of
the conditions nor on the names of their hypotheses. The test of squarefreeness is an example
(Lemma~\ref{lem:Cli.squarefree_spec}): an invariant for each of its three loops
(\lname{JPoints.Cli.EuclidInvariant} for Euclid's algorithm), and the degree of the divisor plus
one as the measure of the last. The lemmas whose titles begin with ``The contract of'' are such
triples; a statement on the value of a component without a loop, like
Lemma~\ref{lem:Printf.emitPoint_eq_render}, is an equation.

\paragraph{Integers and words.} A \texttt{long} of the C code is an integer, and so is a GMP
integer. Where the C code computes in \texttt{long}s, the model takes the integer result, and the
proofs show that the \texttt{long} holds it at the calls: the products of two coordinates of
absolute value at most $\mathtt{MAX\_HEIGHT} = \lfloor \sqrt{2^{63} - 1} \rfloor$ in the Kummer
equation, or the coordinates of a point on the path of machine words of the exact check
(\lname{JPoints.CheckGlobals.abs\_le\_maxHeight\_of\_wordsPath}). The division and the remainder of
the C language truncate on signed operands, and the model uses \lname{Int.tdiv} and
\lname{Int.tmod} for them. Where the 64-bit arithmetic of the C code decides something (the words
of the bit arrays of the sieve, their masks and shifts, the conversion of a coordinate to an
\texttt{unsigned long}), the model computes in \lname{UInt64} and \lname{Int64}, with the semantics
of the C language on x86-64 (Chapter~\ref{chap:words}).

\paragraph{Doubles and floats.} A \texttt{double} that decides what is searched (in the tube, the
real region, the intervals of the real roots and some index arithmetic of the sieve) is a structure
of its IEEE~754 bits (\lname{JPoints.Double}), each operation of the C language and of the C library
a function on the bits, rounding to nearest with ties to even. The proofs decode a double into an
extended rational number (Definition~\ref{def:doubleValue}) and prove of the operations what the
theorems need (Chapter~\ref{chap:doubles}). The cost model, which chooses the sieving primes and the
mode of the sieve, computes in Lean's \lname{Float}, in the order of the C code and with the
\texttt{log} and \texttt{pow} of the C library of the platform. No proof reads a \lname{Float}: the
proofs hold for every choice that the cost model can make (Definition~\ref{def:IsChoice}), and what
the program prints does not depend on the choice.

\paragraph{The sequential path.} \jpoints{} can run its units of work on several threads. The model
follows the path of one thread, which runs the units of work in turn, whatever the number of threads
that \texttt{-t} asks for. With threads, \jpoints{} prints the points of the units in the order of
the units; that it prints what one thread prints is not proved, and the differential tests compare
runs with threads too. The management of memory and the instrumentation of the C code are left out too, and the
model is that of the default build (Chapter~\ref{chap:ccode}).

\section{The command line}\label{sec:input}

\jpoints{} takes the coefficients $f_0, f_1, \dots, f_d$ of the curve $y^2 = f(x)$ as its first
argument, constant term first, the height bound~$H$ as its second, and the options after them.
\texttt{main} answers \texttt{-c list} alone with the constants of the cost model, refuses fewer
than two arguments with error~2, and then calls \texttt{read\_input}, which reads the arguments into
its globals in this order, exiting through \texttt{error} at the first failure
(\csource{j-points.c}{410}{536}):
\begin{enumerate}
  \item up to seven coefficients (error~3 for more) and their degree, the zero leading coefficients
    not counted, which must be at least~$5$ (error~5);
  \item whether~$f$ is squarefree, by Euclid's algorithm on~$f$ and~$f'$ over~$\QQ$ (error~9,
    \csource{j-points.c}{604}{637});
  \item the height bound, in $[1, \mathtt{MAX\_HEIGHT}]$ (error~4);
  \item the options from left to right, each recognized by its second character (error~6 for a
    word that is no option, an option without its value or a value that the option refuses, and
    error~10 for an argument of \texttt{-c} that \texttt{set\_costs} does not read);
  \item the settlement of the options: one thread under \texttt{-1}; the numbers of sieving primes
    of the stages at most \texttt{num\_primes}, and each at most those of the later stages that are
    not $-1$; the bound on the fourth coordinate that of \texttt{-w}, or the height bound, and none
    under \texttt{-a} without \texttt{-w}.
\end{enumerate}
The search needs of an accepted command line the facts of Definition~\ref{def:Cli.Input.WF}, which
the contract of \texttt{read\_input} gives (Lemma~\ref{lem:Cli.readInput_spec}).

\begin{definition}[The globals of the command line]\label{def:Cli.Input}
  \lean{JPoints.Cli.Input}
  \leanok
  \lname{JPoints.Cli.Input} is the record of the globals that \texttt{read\_input} sets: the
  coefficients $f_0, \dots, f_6$, zero above the degree, and the degree; the height bound; the
  bound \texttt{dbound} on the fourth coordinate and the flag \texttt{dbounded} that says whether
  there is one; the numbers of primes of the three stages of the sieve, $-1$ for the choice of the
  cost model; the flags of \texttt{-q}, \texttt{-1} and \texttt{-a}; the format of \texttt{-f}; the
  size of the bit array; the number \texttt{num\_primes} of primes of the table that the sieve may
  use; the number of threads; and the constants of the cost model, with the flags of those that
  \texttt{-c} set. Its defaults are the values that these globals have before \texttt{read\_input}
  runs.
\end{definition}

\begin{definition}[The refusals of the command line]\label{def:Cli.Failure}
  \lean{JPoints.Cli.Failure, JPoints.Cli.Failure.message, JPoints.Cli.Failure.status}
  \leanok
  \lname{JPoints.Cli.Failure} lists the errors with which \jpoints{} refuses a command line, each a
  case of \texttt{error} of the C code: error~2 for fewer than two arguments, 3 for more than seven
  coefficients, 4 for a height bound that is not read or lies outside $[1, \mathtt{MAX\_HEIGHT}]$,
  5 for a degree below~$5$, 6 for a wrong syntax of the options, 9 for a polynomial that is not
  squarefree, and 10 for an argument of \texttt{-c} that is not read.
  \lname{JPoints.Cli.Failure.message} is the text that \texttt{error} prints, the message between
  blank lines, or for error~2 the usage after a blank line, which error~6 prints after its message;
  and \lname{JPoints.Cli.Failure.status} is the exit status, the number of the error.
\end{definition}

\begin{definition}[\texttt{read\_input}]\label{def:Cli.readInput}
  \lean{JPoints.Cli.readInput}
  \leanok
  \uses{def:Cli.Input, def:Cli.Failure, def:Double.ofDyadic}
  \lname{JPoints.Cli.readInput} \textit{args} is \texttt{read\_input} on the words \textit{args}
  after the name of the program, of which \texttt{main} makes sure that there are at least two: the
  globals that the command line sets, read in the order of the C code (the coefficients of the
  first word, of degree at least~$5$, with $f_6 = 0$ for a quintic; the test of squarefreeness
  \lname{JPoints.Cli.squarefree}; the height bound of the second word, in $[1,
  \mathtt{MAX\_HEIGHT}]$; the options and their settlement), or the first error met.
\end{definition}

\begin{definition}[The globals of an accepted command line]\label{def:Cli.Input.WF}
  \lean{JPoints.Cli.Input.WF}
  \leanok
  \uses{def:poly, def:Cli.Input}
  \lname{JPoints.Cli.Input.WF} \textit{input} says that the globals \textit{input} are as
  \texttt{read\_input} leaves them when it accepts the command line: seven coefficients of a
  squarefree polynomial (Definition~\ref{def:poly}) whose degree, at least~$5$, is that of their
  list; $1 \le H \le \mathtt{MAX\_HEIGHT}$; $1 \le \mathtt{dbound} < 2^{63}$, and the fourth
  coordinate bounded unless \texttt{-a} is given; $1 \le \mathtt{num\_primes} \le
  \mathtt{NUM\_PRIMES}$; the numbers of sieving primes of the stages $-1$ or in $[0,
  \mathtt{num\_primes}]$, each at most those of the later stages that are not~$-1$; and a size of
  the bit array of at least~$1$.
\end{definition}

\begin{lemma}[The test of squarefreeness]\label{lem:Cli.squarefree_spec}
  \lean{JPoints.Cli.squarefree_spec}
  \leanok
  \uses{def:poly}
  Let the array \textit{coeffs} hold the coefficients of a polynomial of degree~$d$, constant term
  first, with $0 < d < 7$: the entries above the index~$d$ are zero and the entry at~$d$ is not.
  If \lname{JPoints.Cli.squarefree} \textit{coeffs} $d$ returns true, then the polynomial of the
  coefficients is squarefree.
\end{lemma}

\begin{proof}
  \leanok
  The arrays of the algorithm hold polynomials over~$\QQ$ of the degrees that it keeps, and each
  remainder differs from its dividend by a multiple of the divisor, so that the two polynomials of
  the algorithm are coprime only if~$f$ and~$f'$ are. When the test passes, the gcd that the
  algorithm leaves has degree~$0$; then~$f$ and~$f'$ are coprime, so that~$f$ is separable, hence
  squarefree.
\end{proof}

\begin{lemma}[The contract of \texttt{read\_input}]\label{lem:Cli.readInput_spec}
  \lean{JPoints.Cli.readInput_spec, JPoints.Cli.Input.WF.of_readInput_eq_ok}
  \leanok
  \uses{def:Cli.Input, def:Cli.Failure, def:Cli.readInput, def:Cli.Input.WF}
  If \lname{JPoints.Cli.readInput} \textit{args} returns the globals \textit{input}, then
  \textit{args} has at least two words and is not \texttt{-c list}, and
  \lname{JPoints.Cli.Input.WF} \textit{input} holds.
\end{lemma}

\begin{proof}
  \leanok
  \uses{def:poly, lem:Cli.squarefree_spec, def:Double.ofDyadic}
  The second word gives the height bound, which neither the empty word, read for a missing second
  word, nor \texttt{list} does. The coefficients lie in an array of seven entries, zero above the
  degree, and their polynomial is squarefree by Lemma~\ref{lem:Cli.squarefree_spec}. The loop over
  the options keeps the facts of the coefficients and of the height bound, \texttt{num\_primes} in
  $[1, \mathtt{NUM\_PRIMES}]$, the numbers of sieving primes at least~$-1$ and the size of the bit
  array at least~$1$, option by option, and the settlement after it gives the rest.
\end{proof}

\formalnote{The contract speaks only of the command lines that \texttt{read\_input} accepts: that
it refuses only what it should, for instance only polynomials that are not squarefree, is not
proved, and the main theorems start from an accepted command line (Chapter~\ref{chap:theorems}).
The words of the command line are Lean strings, in which the runtime has replaced the bytes that are
not valid UTF-8 by the character U+FFFD, where \jpoints{} reads the bytes as they are.}

\section{The curve searched}\label{sec:transform}

The search runs on a curve of its own (\csource{j-points.c}{582}{593},
\csource{j-points.c}{594}{600}). A sextic with $f_0 = 0$ and $f_1 = \pm 1$ is reversed by $x
\mapsto 1/x$ into the quintic $y^2 = x^6 f(1/x)$ with leading coefficient~$f_1$; and a quintic,
given or so obtained, with leading coefficient~$-1$ is negated by $x \mapsto -x$ into the monic
quintic $y^2 = f(-x)$. For a monic quintic the first coordinate of a point is a square
(Lemma~\ref{lem:isSquare_of_mem_kummerPoints}), which the main search uses. The program finds the
points of the curve searched and prints them taken back to the curve given. The specification
makes the same transformation (Definition~\ref{def:searchedCoeffs}) and takes the points back by
Definition~\ref{def:backPoint}, and the points so taken back are exactly the points of the curve
given (Lemma~\ref{lem:isPoint_backPoint_iff}). This rests on the behavior of the Kummer equation
and of the lifting test under the two transformations (Lemma~\ref{lem:lifts_iff_of_reversed}): the
reversal exchanges the first and the third coordinate, and the negation negates the second.

\begin{definition}[The curve searched in the program]\label{def:normalizeCurve}
  \lean{JPoints.normalizeCurve, JPoints.searchedCurve}
  \leanok
  \uses{def:poly, def:searchedCoeffs}
  \lname{JPoints.normalizeCurve} \textit{coeffs} \textit{deg} \textit{quiet} is
  \texttt{normalise\_curve}: from the seven coefficients and the degree of the curve given, the
  curve searched (\lname{JPoints.SearchCurve}: its coefficients, its degree and the flags
  \texttt{reversed} and \texttt{negated}) and the message that names it, empty when the curve is
  not transformed or \textit{quiet} is set. A sextic with $f_0 = 0$ and $|f_1| = 1$ is reversed,
  $f_i$ and $f_{6-i}$ exchanged for $i < 3$, into a quintic; then a quintic with $f_5 = -1$ is
  negated, its odd coefficients negated. \lname{JPoints.searchedCurve} \textit{cs} is the curve
  searched as the specification defines it: the coefficients and the flags of
  \lname{JPoints.searchedCoeffs} \textit{cs}, and the degree of these coefficients.
\end{definition}

\begin{lemma}[The cases of the curve searched]\label{lem:searchedCoeffs_eq}
  \lean{JPoints.searchedCoeffs_eq}
  \leanok
  \uses{def:poly, def:searchedCoeffs}
  Let $f_k$ be the entries of the coefficient list~$cs$, zero past its end. If~$cs$ has degree~$6$,
  $f_0 = 0$ and $|f_1| = 1$, then \lname{JPoints.searchedCoeffs} $cs$ is
  $(f_6, f_5, \dots, f_0)$, reversed and not negated, when $f_1 = 1$, and
  $(f_6, -f_5, f_4, -f_3, f_2, -f_1, f_0)$, reversed and negated, when $f_1 = -1$. Otherwise,
  if~$cs$ has degree~$5$ and $f_5 = -1$, it is $(f_0, -f_1, f_2, -f_3, f_4, -f_5, f_6)$, negated
  and not reversed; and otherwise it is~$cs$ itself, neither reversed nor negated.
\end{lemma}

\begin{proof}
  \leanok
  After a reversal the list has degree~$5$ and the leading coefficient~$f_1$, so that the negation
  follows exactly when $f_1 = -1$; unfolding the definition gives the cases.
\end{proof}

\begin{lemma}[The contract of \texttt{normalise\_curve}]\label{lem:normalizeCurve_spec}
  \lean{JPoints.fst_run_normalizeCurve, JPoints.normalizeCurve_spec}
  \leanok
  \uses{def:poly, def:normalizeCurve}
  Let \textit{coeffs} be seven coefficients whose list has the degree $\mathit{deg} \ge 5$. Then
  \lname{JPoints.normalizeCurve} \textit{coeffs} \textit{deg} \textit{quiet} returns the curve
  \lname{JPoints.searchedCurve} of their list: seven coefficients whose degree, at least~$5$ and
  hence~$5$ or~$6$, is that of their list, of a squarefree polynomial when the polynomial given is
  squarefree. Under \textit{quiet} its message is empty.
\end{lemma}

\begin{proof}
  \leanok
  \uses{def:searchedCoeffs, lem:squarefree_poly_searchedCoeffs, lem:searchedCoeffs_eq}
  For seven coefficients the conditions of the two transformations of \texttt{normalise\_curve}
  are those of Lemma~\ref{lem:searchedCoeffs_eq}, its swaps reverse the coefficients and its
  negations negate the odd ones. The degree is~$5$ after a reversal and that of the curve given
  otherwise, and the two transformations keep the polynomial squarefree
  (Lemma~\ref{lem:squarefree_poly_searchedCoeffs}).
\end{proof}

\begin{lemma}[The Kummer surface under the reversal and the negation]\label{lem:lifts_iff_of_reversed}
  \lean{JPoints.kummer_eq_of_reversed, JPoints.lifts_iff_of_reversed, JPoints.kummer_eq_of_negated, JPoints.lifts_iff_of_negated}
  \leanok
  \uses{def:kummer, def:Lifts}
  Let $f, g \colon \NN \to \ZZ$ be the coefficients of two curves, and~$K_f$, $K_g$ their Kummer
  equations (Definition~\ref{def:kummer}).
  \begin{enumerate}
    \item If $g_k = f_{6-k}$ for $k \le 6$, so that $g(x) = x^6 f(1/x)$, then $K_g(a, b, c, d) =
      K_f(c, b, a, d)$. If moreover $f_5 \ne 0$ or $f_6 \ne 0$, and $f_0 \ne 0$ or $f_1 \ne 0$, so
      that both curves have degree at least~$5$, then for $(a, b, c) \ne 0$ with $K_f(c, b, a, d)
      = 0$, the point $(a, b, c, d)$ lifts on~$g$ exactly when $(c, b, a, d)$ lifts on~$f$
      (Definition~\ref{def:Lifts}).
    \item If $g_k = (-1)^k f_k$ for $k \le 6$, so that $g(x) = f(-x)$, then $K_g(a, b, c, d) =
      K_f(a, -b, c, d)$, and $(a, b, c, d)$ lifts on~$g$ exactly when $(a, -b, c, d)$ lifts
      on~$f$.
  \end{enumerate}
\end{lemma}

\begin{proof}
  \leanok
  \uses{def:auxA2}
  The Kummer equations are identities of polynomials. The auxiliary polynomials of the lifting test
  (Definition~\ref{def:auxA2}) and the sextic form of~$g$ at $(a, b, c, d)$ are those of~$f$ at
  $(a, -b, c, d)$, in every chart. The reversal exchanges~$a$ and~$c$, and so compares different
  charts. For $a, c \ne 0$, square roots~$\alpha$ and~$\beta$ of~$A_2$ and~$B_2$ with $C =
  2\alpha\beta$ give square roots of $a^4 A_2'$ and $a^2 B_2'$ for the reversed curve, by
  polynomial identities that transform the coefficients $v_1, v_0$ of the Mumford representation
  linearly. For $c = 0$ the value $a^2 F(b, a)$ of the sextic form is the square $(ab\alpha +
  a^2\beta)^2$. For $a = 0$ the square~$A_2'$ of the reversed curve comes from~$f_6$ or, for $f_6 =
  0$, from $F(c, b)$; this needs the degree at least~$5$.
\end{proof}

\section{The text of a point}\label{sec:printf}

\jpoints{} prints a point whose coordinates are \texttt{long}s by \texttt{emit\_point}
(\csource{j-points.c}{2085}{2090}), with \texttt{printf} under the format of \texttt{-f}, after it
has taken the point back to the curve given: $a$ and~$c$ exchanged after a reversal, $b$ negated
after a negation, and all four coordinates negated when the first nonzero of the first three is
negative. A point whose coordinates need not be \texttt{long}s it prints by \texttt{printf\_mpz}
with \texttt{gmp\_printf}, after a scan of the format has made each conversion
\texttt{\%[-+~\#\textquotesingle 0-9.]*ld} a \texttt{\%\dots Zd}
(\csource{j-points.c}{2042}{2049}).

The C language defines \texttt{printf} only for a format whose conversions each take a
\texttt{long}. The model admits the formats whose conversions are all
\texttt{\%[\textit{flags}][\textit{width}][.\textit{precision}]ld}, at most four of them, and
\texttt{\%\%}: on these \texttt{printf} of glibc reads them as it documents them, the scan of
\texttt{printf\_mpz} makes exactly their \texttt{l}s \texttt{Z}s, and \texttt{gmp\_printf} then
reads the same pieces (Lemma~\ref{lem:Printf.rewriteFormat_spec}). The program refuses the other
formats (Definition~\ref{def:Program.run}). On the admitted formats \texttt{printf} of glibc and
\texttt{gmp\_printf} of GMP print alike but for four departures of GMP from the C language: of the
flags \texttt{+} and space the last one counts, a \texttt{.} without digits is no precision, the
flag~\texttt{0} pads with zeros also with a precision, and under the flags \texttt{-} and~\texttt{0}
together the padding after the digits is of zeros. So a point is printed in one way when its
coordinates are \texttt{long}s and in another when they are not, and the model follows both.

\begin{definition}[The formats of the points]\label{def:Printf.parseFormat}
  \lean{JPoints.Printf.Conversion, JPoints.Printf.Piece, JPoints.Printf.parseFormat, JPoints.Printf.render}
  \leanok
  A piece of a format (\lname{JPoints.Printf.Piece}) is a character that the format prints as it
  is, or a conversion (\lname{JPoints.Printf.Conversion})
  \texttt{\%[\textit{flags}][\textit{width}][.\textit{precision}]ld}, or the same with~\texttt{Z}
  for~\texttt{l}: its flags among \texttt{-}, \texttt{+}, space, \texttt{\#},
  \texttt{\textquotesingle} and~\texttt{0}, its width, and its precision as the format writes it
  (none, a \texttt{.} without digits, or a \texttt{.} with digits). The pieces
  \lname{JPoints.Printf.parseFormat} $\ell$ $s$ of the format~$s$, for the length modifier~$\ell$,
  \texttt{l} or~\texttt{Z}, are defined when~$s$ is admitted: its conversions all of this form,
  taking the arguments in order, with widths and precisions of at most
  \lname{JPoints.Printf.maxField}, and at most four of them, the four arguments that the program
  passes. The text \lname{JPoints.Printf.render} \textit{conv} $F$ $\mathit{xs}$ of the pieces~$F$
  with the arguments~$\mathit{xs}$ prints each conversion with the next argument by \textit{conv}:
  by \lname{JPoints.Printf.printLong} as \texttt{printf} of glibc prints a \texttt{long}, by
  \lname{JPoints.Printf.printMpz} as \texttt{gmp\_printf} of GMP prints an \texttt{mpz\_t}.
\end{definition}

\begin{definition}[Printing a point]\label{def:Printf.emitPoint}
  \lean{JPoints.Printf.emitPoint, JPoints.Printf.printfMpz, JPoints.Printf.rewriteFormat}
  \leanok
  \uses{def:Printf.parseFormat}
  \lname{JPoints.Printf.emitPoint} \textit{rev} \textit{neg} \textit{format} $a$ $b$ $c$ $d$ is what
  \texttt{emit\_point} prints for the point $(a, b, c, d)$ of the curve searched with \texttt{long}
  coordinates: the point taken back to the curve given ($a$ and~$c$ exchanged under \textit{rev},
  $b$ negated under \textit{neg}, then all four negated when the first nonzero of $a, b, c$ is
  negative) under the pieces of the format for~\texttt{l}, by \lname{JPoints.Printf.printLong}.
  \lname{JPoints.Printf.printfMpz} is what \texttt{printf\_mpz} prints for a point with GMP
  coordinates: the format rewritten by \lname{JPoints.Printf.rewriteFormat}, which makes
  the~\texttt{l} of each conversion \texttt{\%[-+~\#\textquotesingle 0-9.]*ld} a~\texttt{Z}; the
  point taken back in the same way, but negated only after a reversal or a negation; and the point
  under the pieces of the rewritten format for~\texttt{Z}, by \lname{JPoints.Printf.printMpz}.
\end{definition}

\formalnote{For a format that is not admitted the model prints nothing, where the C code passes
the format to \texttt{printf} or \texttt{gmp\_printf} whatever it is; the program refuses such a
format before it prints anything else.}

\begin{lemma}[The format of \texttt{gmp\_printf}]\label{lem:Printf.rewriteFormat_spec}
  \lean{JPoints.Printf.rewriteFormat_spec, JPoints.Printf.parseFormat_run_rewriteFormat}
  \leanok
  \uses{def:Printf.parseFormat, def:Printf.emitPoint}
  If \texttt{printf} reads the characters of \textit{format} with the length modifier~\texttt{l}
  as the pieces~$F$, then \texttt{gmp\_printf} reads those of \lname{JPoints.Printf.rewriteFormat}
  \textit{format} with the length modifier~\texttt{Z} as the same pieces~$F$. In particular, when
  \textit{format} has the pieces~$F$ for~\texttt{l} by \lname{JPoints.Printf.parseFormat}, the
  rewritten format has the pieces~$F$ for~\texttt{Z}.
\end{lemma}

\begin{proof}
  \leanok
  A conversion of an admitted format is its~\texttt{\%}, its flags, width and precision, all of
  them characters that the scan passes over, then its~\texttt{l} and its~\texttt{d}; and the reading
  of a format takes these characters alike for~\texttt{l} and for~\texttt{Z}. The invariant of the
  outer loop of the scan splits the copy and the format at the index of the loop into a prefix,
  read alike with~\texttt{Z} and with~\texttt{l} into the same state and pieces, and the same rest;
  the state is outside a conversion, or after the~\texttt{l} of one whose~\texttt{d} comes next.
\end{proof}

\begin{lemma}[The text of a printed point]\label{lem:Printf.emitPoint_eq_render}
  \lean{JPoints.Printf.emitPoint_eq_render, JPoints.Printf.run_printfMpz_eq_render}
  \leanok
  \uses{def:backPoint, def:Printf.parseFormat, def:Printf.emitPoint}
  Let~$F$ be the pieces \lname{JPoints.Printf.parseFormat} of \textit{format} for~\texttt{l}, and
  let the map back (Definition~\ref{def:backPoint}) under \textit{rev} and \textit{neg} take $(a,
  b, c, d)$ to $(a', b', c', d')$. Then
  \[
    \mathrm{emitPoint}\ \mathit{rev}\ \mathit{neg}\ \mathit{format}\ a\ b\ c\ d
      = \mathrm{render}\ \mathrm{printLong}\ F\ [a', b', c', d']\,,
  \]
  and, when the first nonzero of $a, b, c$ is positive, \lname{JPoints.Printf.printfMpz} prints the
  same with \lname{JPoints.Printf.printMpz} in place of \lname{JPoints.Printf.printLong}.
\end{lemma}

\begin{proof}
  \leanok
  \uses{lem:Printf.rewriteFormat_spec}
  The first is the definition: \texttt{emit\_point} takes the point back as the specification
  does. For the second, \texttt{gmp\_printf} reads the rewritten format with the pieces~$F$
  (Lemma~\ref{lem:Printf.rewriteFormat_spec}), and without a reversal or a negation
  \texttt{printf\_mpz} does not change the signs, which a point with its first nonzero coordinate
  positive does not need.
\end{proof}

\section{The exact check}\label{sec:check}

The sieve hands the exact check the triples $(a, b, c)$ that it keeps, coprime, with $a \ge 0$ and
the first nonzero coordinate positive. \texttt{check\_one\_point}
(\csource{j-points.c}{2180}{2257}) computes the Kummer equation $K_2 d^2 + K_1 d + K_0 = 0$ over the
triple, with~$K_0$ from the fifteen coefficients of the quartic that \texttt{kummer\_init} computes
once (\csource{j-points.c}{1421}{1463}), and solves it by its discriminant, in the order of the
roots of the specification (Definition~\ref{def:kummerRoots}). \texttt{check\_one\_point\_final}
(\csource{j-points.c}{2097}{2178}) writes each rational root as $u/g$ in lowest terms with $g > 0$,
which gives the candidate $(ga, gb, gc, u)$, and decides by the size of~$g$ and~$u$ where to check
it: on machine words when its coordinates are at most \texttt{MAX\_HEIGHT} in absolute value, so
that the products of the lifting test (\texttt{check\_lifts}, \csource{j-points.c}{1860}{1947}) fit
a \texttt{long}, and otherwise on GMP integers, when the candidate keeps the bounds of the search.
On either path a point that lifts is printed (Section~\ref{sec:printf}) and counted.

\begin{definition}[The Kummer equation of the program]\label{def:kummerInit}
  \lean{JPoints.kummerInit, JPoints.kummerEqn}
  \leanok
  \lname{JPoints.kummerInit} \textit{coeffs} is the record \lname{JPoints.KummerCoeffs} of the
  fifteen coefficients $k_{400}, k_{310}, \dots, k_{004}$ of the quartic~$K_0$, $k_{ijk}$ the
  coefficient of $a^i b^j c^k$, each computed from the coefficients of the curve as
  \texttt{kummer\_init} computes it, from $k_{400} = f_1^2 - 4f_0f_2$ to $k_{004} = f_5^2 -
  4f_4f_6$. \lname{JPoints.kummerEqn} $k$ \textit{coeffs} $a$ $b$ $c$ is the record
  \lname{JPoints.KummerEqn} of the coefficients $K_0$, $K_1$, $K_2$ of the Kummer equation over
  $(a, b, c)$ as \texttt{kummer\_eqn} computes them from the products $a^2, b^2, c^2, ab, ac, bc$:
  $K_0$ from the coefficients~$k$, and
  \[
    K_1 = -4a\,(f_0a^2 + f_2ac) - 2b\,(f_1a^2 + f_3ac + f_5c^2) - 4c\,(f_4ac + f_6c^2)\,, \qquad
    K_2 = b^2 - 4ac\,.
  \]
\end{definition}

\begin{lemma}[The Kummer equation of the program is that of the specification]\label{lem:kummerEqn_kummerInit}
  \lean{JPoints.kummerEqn_kummerInit}
  \leanok
  \uses{def:kummer, def:kummerInit}
  For the coefficients \textit{coeffs} and a triple $(a, b, c)$, the coefficients of
  \lname{JPoints.kummerEqn} (\lname{JPoints.kummerInit} \textit{coeffs}) \textit{coeffs} $a$ $b$ $c$
  are $K_0(a, b, c)$, $K_1(a, b, c)$ and $K_2(a, b, c)$ of the specification
  (Definition~\ref{def:kummer}) for the polynomial with the coefficients \textit{coeffs}, zero past
  their end.
\end{lemma}

\begin{proof}
  \leanok
  The three equations are identities of polynomials in the coefficients and the coordinates.
\end{proof}

\begin{definition}[The lifting test of the program]\label{def:checkLifts}
  \lean{JPoints.checkLifts, JPoints.checkLiftsMpz}
  \leanok
  \lname{JPoints.checkLifts} \textit{coeffs} $a$ $b$ $c$ $d$ is the lifting test
  \texttt{check\_lifts} on the point $(a : b : c : d)$ of the Kummer surface of the curve searched,
  given by machine words: the origin $(0 : 0 : 0 : d)$ passes; $(0 : 0 : c : d)$ passes when~$f_6$
  is a square; $(0 : b : c : d)$ with $b \ne 0$ when~$f_6$ and $b^6 f(c/b)$ are squares; and
  $(a : b : c : d)$ with $a \ne 0$ when~$A_2$ is a square and, if $A_2 = 0$, when~$B_2$ is one,
  for the auxiliary polynomials~$A_2$ and~$B_2$ of the test (Definition~\ref{def:auxA2}); $0$ counts
  as a square. \lname{JPoints.checkLiftsMpz} is the same test, \texttt{check\_lifts\_mpz}, on GMP
  integers, with the products in the order of that code.
\end{definition}

\begin{lemma}[The contract of the lifting test]\label{lem:checkLifts_spec}
  \lean{JPoints.checkLifts_spec, JPoints.checkLiftsMpz_spec}
  \leanok
  \uses{def:BruteForce.lifts, def:checkLifts}
  Let $(a, b, c) \ne 0$. Then \lname{JPoints.checkLiftsMpz} \textit{coeffs} $a$ $b$ $c$ $d$
  returns \lname{JPoints.BruteForce.lifts} \textit{coeffs} $a$ $b$ $c$ $d$
  (Definition~\ref{def:BruteForce.lifts}), and so does \lname{JPoints.checkLifts} when moreover $0
  \le b < 2^{64}$ in case $a = 0$.
\end{lemma}

\begin{proof}
  \leanok
  \uses{def:auxA2, lem:BruteForce.lifts_eq_ite}
  The loops compute $b^6 f(c/b)$ by Horner's rule, which named invariants record; the values that
  the two tests compute for~$A_2$ and~$B_2$ are those of Definition~\ref{def:auxA2}; the test of
  GMP for squares decides the squares; and on machine words the conversion of~$b$ to an
  \texttt{unsigned long} is the identity for $0 \le b < 2^{64}$. Then the case tree is that of the
  brute-force test (Lemma~\ref{lem:BruteForce.lifts_eq_ite}).
\end{proof}

\begin{definition}[The globals of the exact check and the output]\label{def:CheckGlobals}
  \lean{JPoints.CheckGlobals, JPoints.Output, JPoints.Output.print}
  \leanok
  \lname{JPoints.CheckGlobals} is what the exact check reads of the globals: the curve searched
  (\lname{JPoints.SearchCurve}: its coefficients, its degree and the flags \texttt{reversed} and
  \texttt{negated}), the coefficients of the quartic~$K_0$ (\lname{JPoints.KummerCoeffs}), the
  height bound, the bound \texttt{dbound} on the fourth coordinate and the flag \texttt{dbounded}
  that says whether there is one, the flags of \texttt{-a} and \texttt{-1}, and the format of
  \texttt{-f}. \lname{JPoints.Output} is the output of the search: the text printed so far and the
  count \texttt{total} of the points found. \lname{JPoints.Output.print} \textit{out} $s$ appends
  the text~$s$ to the text of \textit{out}.
\end{definition}

\begin{definition}[The globals of the exact check as \texttt{main} sets them]\label{def:CheckGlobals.WF}
  \lean{JPoints.CheckGlobals.WF}
  \leanok
  \uses{def:poly, def:CheckGlobals, def:kummerInit}
  \lname{JPoints.CheckGlobals.WF} \textit{gl} says that the globals of the exact check are as
  \texttt{main} sets them from a command line that \texttt{read\_input} accepts: the curve searched
  has seven coefficients, whose degree, at least~$5$ and hence~$5$ or~$6$, is that of their list
  (Definition~\ref{def:poly}); the coefficients of the quartic are \lname{JPoints.kummerInit} of
  them; $1 \le H \le \mathtt{MAX\_HEIGHT}$; the bound on the fourth coordinate is at least~$1$ when
  there is one; and there is one unless \texttt{-a} is given.
\end{definition}

The bounds of the specification come from the globals \textit{gl}: the height bound~$H$, as a
natural number (\lname{JPoints.CheckGlobals.H}), the flag \textit{all} of \texttt{-a}, and the
bound~$W$ on the fourth coordinate if there is one, none otherwise (\lname{JPoints.CheckGlobals.W}).
Within these bounds the points of the curve searched over a triple~$t$, in the order of the roots
of the Kummer equation (Definition~\ref{def:HasPointOver}), are \lname{JPoints.CheckGlobals.points}
\textit{gl} $t$. Let~$H'$ be the height bound, or \texttt{MAX\_HEIGHT} under \texttt{-a},
and~$d_{\max}$ the bound on the fourth coordinate when there is one below \texttt{MAX\_HEIGHT},
and \texttt{MAX\_HEIGHT} otherwise. A point of the curve searched lies on the path of machine words
(\lname{JPoints.CheckGlobals.wordsPath}) when its first three coordinates are at most~$H'$ and its
fourth at most~$d_{\max}$ in absolute value. Its text (\lname{JPoints.CheckGlobals.text}) is that
of \lname{JPoints.Printf.emitPoint} on this path and that of \lname{JPoints.Printf.printfMpz}
otherwise, and the output with the points~$L$ appended is the output with their texts after its
text and their number added to its count (\lname{JPoints.Output.appendPoints}).

\begin{definition}[The exact check]\label{def:checkOnePoint}
  \lean{JPoints.checkOnePointFinal, JPoints.checkOnePoint}
  \leanok
  \uses{def:CheckGlobals, def:kummerInit, def:Printf.emitPoint, def:checkLifts}
  The check \texttt{check\_one\_point\_final} of the root $d_1/d_2$ of the Kummer equation over
  $(a, b, c)$ is \lname{JPoints.checkOnePointFinal} \textit{gl} \textit{out} $a$ $b$ $c$ $d_1$
  $d_2$. Let $u/v = d_1/d_2$ in lowest terms and $m = \max(a, |b|, |c|)$; let~$H'$ be the height
  bound, or \texttt{MAX\_HEIGHT} under \texttt{-a}, and $h = \lfloor H'/m \rfloor$ the bound on the
  scale; and let~$d_{\max}$ be the bound on the fourth coordinate when there is one below
  \texttt{MAX\_HEIGHT}, and \texttt{MAX\_HEIGHT} otherwise.
  \begin{enumerate}
    \item When $|u| \le d_{\max}$ and $|v| \le h$, it checks on machine words the point $(ga,
      gb, gc, \pm u)$, where $g = |v|$ and the sign is that of~$v$: the lifting test by
      \lname{JPoints.checkLifts}, the printing by \lname{JPoints.Printf.emitPoint}.
    \item Otherwise, when $|u| \le \mathtt{dbound}$ or the fourth coordinate is not bounded, and
      $|v| \le h$ or \texttt{-a} is given, it checks the same on GMP integers: the lifting test
      by \lname{JPoints.checkLiftsMpz}, the printing by \lname{JPoints.Printf.printfMpz}.
    \item Otherwise it checks nothing.
  \end{enumerate}
  A point that passes is appended to the output and counted, and the result says whether one was.
  \lname{JPoints.checkOnePoint} \textit{gl} \textit{out} $a$ $b$ $c$ is \texttt{check\_one\_point}:
  it computes the Kummer equation over $(a, b, c)$ by \lname{JPoints.kummerEqn} from the
  coefficients of \textit{gl}, and its discriminant $\Delta = K_1^2 - 4K_0K_2$. For $\Delta = 0$ it
  checks the root $-K_1/(2K_2)$, unless $K_2 = K_1 = 0$; for~$\Delta$ a nonzero square it checks
  the root $-K_0/K_1$ when $K_2 = 0$, and otherwise $(\sqrt\Delta - K_1)/(2K_2)$ and then
  $-(\sqrt\Delta + K_1)/(2K_2)$, the second not under \texttt{-1} when the first gave a point;
  for~$\Delta$ not a square nothing. Each root is checked by \lname{JPoints.checkOnePointFinal}.
\end{definition}

\formalnote{The bound $h = H'/m$ and the comparison of \texttt{dbound} with \texttt{MAX\_HEIGHT}
are made in \texttt{unsigned long}s in the C code, the type of \texttt{MAX\_HEIGHT}; the model
makes them in the integers, where they agree for the positive values of the calls. Under
\texttt{-a} a candidate whose scale exceeds~$h$ goes to the GMP integers with no bound on the scale
at all, as in the C code.}

\begin{lemma}[The candidate of a root]\label{lem:isPoint_pointOfRoot_iff}
  \lean{JPoints.pointOfRoot, JPoints.isPoint_pointOfRoot_iff}
  \leanok
  \uses{def:kummer, def:InBounds, def:IsPoint, def:BruteForce.lifts}
  For a triple $(a, b, c)$ and a rational number $x = u/g$ in lowest terms with $g > 0$, the
  candidate \lname{JPoints.pointOfRoot} $a$ $b$ $c$ $x$ is $(ga, gb, gc, u)$. Let $\gcd(a, b, c) =
  1$, the first nonzero of $a, b, c$ positive, let~$x$ be a root of $K_2x^2 + K_1x + K_0$ over $(a,
  b, c)$ for the coefficient list~$cs$, and let the candidate lie within the bounds~$H$,
  \textit{all} and~$W$ of the search (Definition~\ref{def:InBounds}). Then the candidate is a point,
  $\mathrm{IsPoint}\ cs\ H\ \mathit{all}\ W$ (Definition~\ref{def:IsPoint}), exactly when
  \lname{JPoints.BruteForce.lifts} $cs$ holds for it.
\end{lemma}

\begin{proof}
  \leanok
  \uses{def:auxA2, def:Lifts, def:IsNormalized, lem:lifts_iff_of_kummer_eq_zero}
  The Kummer equation is homogeneous of degree~$4$, so that $K(ga, gb, gc, u) = g^4 (K_2x^2 + K_1x
  + K_0) = 0$: the candidate lies on the Kummer surface. It is normalized
  (Definition~\ref{def:IsNormalized}): $\gcd(g, u) = 1$ and $\gcd(a, b, c) = 1$ make its
  coordinates coprime, and $g > 0$ keeps the sign of the first nonzero one. On the Kummer surface
  the test of the brute force decides $\mathrm{Lifts}$
  (Lemma~\ref{lem:lifts_iff_of_kummer_eq_zero}).
\end{proof}

\begin{lemma}[The contract of \texttt{check\_one\_point\_final}]\label{lem:checkOnePointFinal_spec}
  \lean{JPoints.checkOnePointFinal_spec}
  \leanok
  \uses{def:kummer, def:IsPoint, def:orderedTriples, def:CheckGlobals, def:CheckGlobals.WF, def:checkOnePoint, lem:isPoint_pointOfRoot_iff, def:Printf.emitPoint}
  Let \lname{JPoints.CheckGlobals.WF} \textit{gl}, let $(a, b, c)$ be a triple of
  $\mathrm{boxTriples}\ H$ (Definition~\ref{def:orderedTriples}), let $d_2 \ne 0$, and let $x =
  d_1/d_2$ be a root of the Kummer equation over $(a, b, c)$ of the curve searched. Then
  \lname{JPoints.checkOnePointFinal} \textit{gl} \textit{out} $a$ $b$ $c$ $d_1$ $d_2$ appends to
  \textit{out} the candidate \lname{JPoints.pointOfRoot} $a$ $b$ $c$ $x$ exactly when it is a point
  ($\mathrm{IsPoint}$ for the curve searched and the bounds of \textit{gl}), with its text and its
  count, and returns whether it is one.
\end{lemma}

\begin{proof}
  \leanok
  \uses{def:Lifts, def:IsNormalized, def:InBounds, lem:mem_orderedTriples_iff, def:BruteForce.lifts, def:checkLifts, lem:checkLifts_spec}
  The triple of the box is coprime, with $a \ge 0$ and its first nonzero coordinate positive
  (Lemma~\ref{lem:mem_orderedTriples_iff}), so that the candidate is normalized. On the path of
  machine words the bounds of the search hold (Definition~\ref{def:InBounds}): $g \le h$ means $gm
  \le H'$. On the path of GMP integers the check tests exactly whether they hold, and otherwise the
  candidate is no point. On either path the lifting test is that of the brute force
  (Lemma~\ref{lem:checkLifts_spec}; on machine words $b \ge 0$ when $a = 0$), which decides
  $\mathrm{Lifts}$ on the candidate (Lemma~\ref{lem:isPoint_pointOfRoot_iff}).
\end{proof}

\begin{lemma}[The contract of \texttt{check\_one\_point}]\label{lem:checkOnePoint_spec}
  \lean{JPoints.checkOnePoint_spec}
  \leanok
  \uses{def:orderedTriples, def:firstIf, def:HasPointOver, def:CheckGlobals, def:CheckGlobals.WF, def:checkOnePoint, def:Printf.emitPoint}
  Let \lname{JPoints.CheckGlobals.WF} \textit{gl} and $(a, b, c) \in \mathrm{boxTriples}\ H$, and
  let \textit{one} be the flag of \texttt{-1}. Then \lname{JPoints.checkOnePoint} \textit{gl}
  \textit{out} $a$ $b$ $c$ appends to \textit{out} the points
  \[
    \mathrm{firstIf}\ \mathit{one}\ (\mathit{gl}.\mathrm{points}\ (a, b, c))
  \]
  (Definition~\ref{def:firstIf}), the points over $(a, b, c)$ (Definition~\ref{def:HasPointOver}),
  in the order of the roots of the Kummer equation, only the first under \texttt{-1}. Under
  \texttt{-1} it returns true exactly when a point lies over $(a, b, c)$.
\end{lemma}

\begin{proof}
  \leanok
  \uses{def:kummer, def:IsPoint, def:kummerRoots, lem:mem_kummerRoots_iff, def:BruteForce.lifts, def:kummerInit, lem:kummerEqn_kummerInit, lem:isPoint_pointOfRoot_iff, lem:checkOnePointFinal_spec}
  The check computes the Kummer equation of the specification
  (Lemma~\ref{lem:kummerEqn_kummerInit}), and the case tree of the discriminant gives the roots of
  Definition~\ref{def:kummerRoots} in their order (Lemma~\ref{lem:mem_kummerRoots_iff}), with the
  linear root when $K_2 = 0$, so that no squarefreeness of the curve is needed. Each root is checked
  by Lemma~\ref{lem:checkOnePointFinal_spec}, whose candidates are those over the triple
  (Lemma~\ref{lem:isPoint_pointOfRoot_iff}); under \texttt{-1} the second root is checked only when
  the first gave no point, which leaves the first point.
\end{proof}

\section{The program}\label{sec:program}

\texttt{main} reads the command line (Section~\ref{sec:input}) and searches
(\csource{j-points.c}{359}{372}): it transforms the curve (Section~\ref{sec:transform}), analyzes
the plane of $(b, c)$ (\texttt{tube\_init}, Chapter~\ref{chap:tube}), computes the condition at~$2$
(\texttt{twoadic\_init}, Chapter~\ref{chap:sieve-numbers}) and calls \texttt{find\_points}, which
chooses the sieve and then checks the points over $(0, 0, 1)$ (\csource{j-points.c}{1304}{1310}),
the double points, and the units of work of the main search, each a first coordinate~$a$ whose rows
$(a, b)$ \texttt{run\_a} sieves (Chapter~\ref{chap:words}). So the triples come in the order of the
specification: first those on the conic $b^2 = 4ac$, that is $(0, 0, 1)$ and those of the double
points, then the others, unit by unit (Lemma~\ref{lem:flatMap_orderedTriples}).

\begin{definition}[\texttt{find\_points}]\label{def:findPoints}
  \lean{JPoints.findPoints}
  \leanok
  \uses{def:CheckGlobals, def:Plan, def:runA, def:Cli.Input, def:Double.add, def:Double.ofDyadic, def:SieveTables.WF, def:SiftData.WF, def:checkOnePoint, def:initFmodpsquare, def:message}
  \lname{JPoints.findPoints} $g$ \textit{plan} \textit{sd} $t$ \textit{out} is
  \texttt{find\_points} for the globals~$g$ of the search (\lname{JPoints.SearchGlobals}, the
  globals of the exact check among them), from the initial plan of the sieve
  (Definition~\ref{def:Plan}), the static data \textit{sd} of the sieve and the state~$t$ of the
  analysis of the tube that \texttt{tube\_init} leaves, after the output \textit{out}: the set-up,
  which makes the choice of the sieve and its tables and may end the search; the points over $(0,
  0, 1)$ (\lname{JPoints.findPointsAtInfinity}); the double points (\texttt{find\_double\_points}),
  each ending the search when it finds a point under \texttt{-1}; and the main search over the
  units of work (\lname{JPoints.runUnits}, with \texttt{run\_a} of Definition~\ref{def:runA}) in
  the state of the one thread. It returns the totals of the counters and the output.
\end{definition}

\begin{lemma}[The points over $(0, 0, 1)$]\label{lem:findPointsAtInfinity_spec}
  \lean{JPoints.findPointsAtInfinity_spec}
  \leanok
  \uses{def:firstIf, def:HasPointOver, def:CheckGlobals, def:CheckGlobals.WF, def:Printf.emitPoint, def:checkOnePoint}
  Let \lname{JPoints.CheckGlobals.WF} \textit{gl}, and let \textit{one} be the flag of \texttt{-1}.
  Then the search of \texttt{find\_points} over $(0, 0, 1)$, \lname{JPoints.findPointsAtInfinity}
  \textit{gl} \textit{out}, appends to \textit{out} the points
  \[
    \mathrm{firstIf}\ \mathit{one}\ (\mathit{gl}.\mathrm{points}\ (0, 0, 1))\,,
  \]
  and returns whether \texttt{-1} is given and a point lies over $(0, 0, 1)$.
\end{lemma}

\begin{proof}
  \leanok
  \uses{def:poly, def:kummer, def:auxA2, def:Lifts, def:IsNormalized, def:InBounds, def:IsPoint, def:orderedTriples, def:kummerRoots, lem:mem_orderedTriples_iff, lem:mem_kummerRoots_iff, def:BruteForce.lifts, def:checkOnePoint, lem:isPoint_pointOfRoot_iff, lem:checkOnePointFinal_spec}
  Over $(0, 0, 1)$ the Kummer equation has $K_2 = 0$, $K_1 = -4f_6$ and $K_0 = f_5^2 - 4f_4f_6$. A
  quintic has $K_1 = K_2 = 0$ and no point there, and a point over $(0, 0, 1)$ lifts only when~$f_6$
  is a square. For a sextic whose~$f_6$ is a square the search checks the one root $(f_5^2 -
  4f_4f_6)/(4f_6)$ of Definition~\ref{def:kummerRoots} by \texttt{check\_one\_point\_final}
  (Lemma~\ref{lem:checkOnePointFinal_spec}).
\end{proof}

\begin{definition}[The messages]\label{def:message}
  \lean{JPoints.message}
  \leanok
  \uses{def:Cli.Input, def:Plan}
  \lname{JPoints.message} $g$ $n$ \textit{total} is the text of message~$n$ of \jpoints{} for the
  globals~$g$ (\lname{JPoints.MessageGlobals}): $0$ the banner of the version; $1$ the end of a
  search that a prime of rate~$0$ leaves without points; $2$ the number of the points found; $4$
  the mode of the sieve, the condition at~$2$, the numbers of the primes of the stages and the
  constants that \texttt{-c} set; $5$ the curve given; $6$ the height bound, with the bound on the
  fourth coordinate when it differs from it; $7$ the sieving primes of the stages; $8$ the rates of
  the best prime, of the last of each stage and of the worst; $12$ the numbers of the survivors of
  the three stages of the sieve; and nothing for any other~$n$. The callers leave the messages out
  under \texttt{-q}.
\end{definition}

\begin{definition}[The program]\label{def:Program.run}
  \lean{JPoints.Program.mainSearch, JPoints.Program.run, JPoints.Program.main}
  \leanok
  \uses{def:Cli.Input, def:Cli.Failure, def:Cli.readInput, def:CheckGlobals, def:kummerInit, def:Printf.parseFormat, def:normalizeCurve, def:findPoints, def:message, def:Plan, def:Double.add, def:Double.ofDyadic, def:realRootsInit, def:rowRegion}
  \lname{JPoints.Program.run} \textit{args} is \texttt{main} of \jpoints{} on the words
  \textit{args} after the name of the program: the text that it prints on the standard output and
  its exit status. It makes the tables of \texttt{init\_main}; under \texttt{-c list} alone it
  prints the constants of the cost model, with the status~$0$; with fewer than two words it fails
  with error~2; otherwise it reads the command line (Definition~\ref{def:Cli.readInput}), and on an
  error it prints its message, with its status (Definition~\ref{def:Cli.Failure}); it refuses a
  format of \texttt{-f} that is not admitted (Definition~\ref{def:Printf.parseFormat}) with the
  status~$11$; and then it prints the head of the report, unless \texttt{-q} is given, and the text
  of its search, with the status~$0$. \lname{JPoints.Program.mainSearch} $t$ \textit{input} is that
  search for the tables~$t$ and the globals \textit{input} of the command line: the curve searched
  (Definition~\ref{def:normalizeCurve}), the masks of the ends of a row, the size of the bit array
  in words, the coefficients of the quartic (Definition~\ref{def:kummerInit}), the analysis of the
  plane of $(b, c)$ and its sample, the condition at~$2$ when $H \ge 500$, the search
  (Definition~\ref{def:findPoints}), and unless \texttt{-q} is given messages~$12$ and~$2$
  (Definition~\ref{def:message}). \lname{JPoints.Program.main} is the program
  \texttt{j-points-model}: it writes the text of \lname{JPoints.Program.run} on the standard output
  and exits with its status.
\end{definition}

\formalnote{The refusal of a format of \texttt{-f} that the model does not admit, with the
status~$11$, models no C code: \jpoints{} passes any format to \texttt{printf}. The admitted formats
are those that \texttt{printf} of glibc, the scan of \texttt{printf\_mpz} and \texttt{gmp\_printf}
read as the model does (Section~\ref{sec:printf}). The program of the model computes the whole text
and writes it at the end, where \jpoints{} prints as it goes. When the reader of the output is
gone, the write fails and \texttt{j-points-model} exits quietly with the status~$141$, which the
shell reports for \jpoints{}, ended by the signal \texttt{SIGPIPE}, which the Lean runtime
ignores.}

\begin{definition}[The globals of the search as \texttt{main} sets them]\label{def:IsSearchGlobals}
  \lean{JPoints.IsSearchGlobals}
  \leanok
  \uses{def:poly, def:Cli.Input.WF, def:CheckGlobals, def:CheckGlobals.WF, def:Tables.WF, def:TwoAdic.Sound, def:realRootsInit, def:NegativeIntervals.Sound}
  \lname{JPoints.IsSearchGlobals} $g$ says that the globals~$g$ of the search are as \texttt{main}
  sets them from a command line that \texttt{read\_input} accepts: the globals of the exact check
  satisfy \lname{JPoints.CheckGlobals.WF} (Definition~\ref{def:CheckGlobals.WF}) and have the
  height bound, the bound on the fourth coordinate and its flag of the globals of the command line,
  which satisfy \lname{JPoints.Cli.Input.WF} (Definition~\ref{def:Cli.Input.WF}); the tables of
  \texttt{init\_main} satisfy \lname{JPoints.Tables.WF} (Definition~\ref{def:Tables.WF}); the
  choice reads the curve searched; the coefficients of the analysis of the tube are those of
  \texttt{tube\_init}; the condition at~$2$ admits every point
  (Definition~\ref{def:TwoAdic.Sound}); the intervals are those of \texttt{real\_roots\_init}
  (Definition~\ref{def:realRootsInit}), on which~$f$ is negative
  (Definition~\ref{def:NegativeIntervals.Sound}) unless the sieve has no primes at all
  ($\mathtt{sieve\_primes3} = 0$, which \texttt{-N} with a value of at most~$0$ gives); the bit
  array has at least one word, and the masks of the ends of a row are those of \texttt{main}; the
  polynomial of the curve searched is squarefree (Definition~\ref{def:poly}); and the height bound
  is at most~$2^{26}$.
\end{definition}

The contracts of this chapter meet in the proof of the main theorem. On a command line that
\texttt{read\_input} accepts, its contract gives the facts of Definition~\ref{def:Cli.Input.WF}
(Lemma~\ref{lem:Cli.readInput_spec}). From them the contract of \texttt{normalise\_curve}
(Lemma~\ref{lem:normalizeCurve_spec}) shows that the globals of the exact check that \texttt{main}
sets (\lname{JPoints.checkGlobalsOf}) satisfy Definition~\ref{def:CheckGlobals.WF}, and the
contracts of \texttt{init\_main}, \texttt{tube\_init}, \texttt{twoadic\_init} and
\texttt{real\_roots\_init} give the rest of Definition~\ref{def:IsSearchGlobals}: Sturm's theorem
(Definition~\ref{def:SturmTheorem}) is needed for the intervals unless $\mathtt{sieve\_primes3} =
0$, and the height bound~$2^{26}$, below which the real region keeps every point
(Chapter~\ref{chap:region}), is a hypothesis of the main theorem. The contract of
\texttt{find\_points} (Lemma~\ref{lem:findPoints_spec}) then says that the search prints the text of
the choice of the sieve and the points over the triples of the specification, each by the exact
check of this chapter, and the contract of the search of \texttt{main}
(Lemma~\ref{lem:Program.mainSearch_spec}) gives the main theorem
(Theorem~\ref{thm:Program.run_eq_of_readInput_eq_ok_of_quiet}).
